\chapter{系统实现}
\label{chap:implementation}

\section{驱动函数接口静态分析与识别模块实现}

\subsection{CodeQL 查询实现}

查询流水线依次执行四阶段：MMIO 表达式收集 $\to$ 基址数据流追踪 $\to$ 驱动函数判定 $\to$ 调用关系查询，各阶段输出作为下一阶段输入。

\subsubsection{MMIO 操作定位查询实现}

针对三种 MMIO 定位模式，本研究设计了相应的查询规则。查询的核心设计思想是将"寄存器字段访问"追溯到一个可被 CodeQL 稳定描述的基对象表达式。代码~\ref{lst:mmio-collect}展示了 MMIO 表达式收集查询，通过 \texttt{import libraries.MMIOAnalyzer} 引入 MMIO 分析库，从数据库中提取所有符合条件的 MMIO 表达式及其所属函数、文件路径、行号和类型信息。

\begin{lstlisting}[style=CodeQL,language=CodeQL, caption={MMIO表达式收集查询示例},label=lst:mmio-collect]
import cpp
import libraries.MMIOAnalyzer

from Function func, InterestingMMIOExpr expr
where isInterestingMMIOFunction(func) and
     expr.getEnclosingFunction() = func
select expr.toString(),
       expr.getEnclosingFunction().getName(),
       expr.getFile().getAbsolutePath().toString(),
       expr.getLocation().getStartLine()
\end{lstlisting}

对于"间接初始化"和"基址直传"模式，查询采用 CodeQL 的数据流追踪能力，将外设基址常量（如 \texttt{USART1\_BASE}）作为源点，沿赋值语句和函数参数传播追踪到最终的字段解引用操作（如 \texttt{huart->Instance->DR}），如代码~\ref{lst:taint-tracking}所示。

\begin{lstlisting}[style=CodeQL,language=CodeQL, caption={常量基址数据流追踪查询示例},label=lst:taint-tracking]
predicate isMMIOTracedExpr(Expr e) {
    exists(DataFlow::Node source, DataFlow::Node sink,
           ConstantExpr src_mmio |
        source.asExpr() = src_mmio and
        isConstantBigValueExpr(src_mmio) and
        sink.asExpr() = e and
        TaintFlow::flow(source, sink) and
        not sink.asExpr().getFile() instanceof DriverFile
    )
}
\end{lstlisting}

对于"基址直传"模式，查询通过函数调用关系和参数类型匹配定位调用点：当函数的第 0 个参数类型为 MMIO 结构体指针且实参是经过数据流追踪的常量表达式时，该调用点即为候选的 MMIO 使用点。三种查询的结果经聚合后形成完整的 MMIO 使用链路，为驱动函数识别提供输入。

\subsubsection{驱动函数定位查询实现}

驱动函数定位的判定规则见第~\ref{sec:design-driver-loc}~节，其 CodeQL 实现如代码~\ref{lst:driver-function-rule}所示。MMIO 访问包括直接字段访问（\texttt{MMIOFieldAccess}）和间接字段访问（\texttt{DriverFieldAccess}），两类访问的函数均被识别为驱动函数。

\begin{lstlisting}[style=CodeQL,language=CodeQL, caption={驱动函数判定规则},label=lst:driver-function-rule]
class DriverFunction extends Function
{
    DriverFunction() {
        exists(DriverFieldAccess fa | fa.getEnclosingFunction() = this)
        or
        exists(MMIOFieldAccess fa | fa.getEnclosingFunction() = this)
    }
}
\end{lstlisting}

该查询定义了 \texttt{DriverFunction} 类，该类继承自 \texttt{Function}，通过检查函数内部是否包含 \texttt{DriverFieldAccess} 或 \texttt{MMIOFieldAccess} 来判定函数是否为驱动函数，与设计章中"凡是包含 MMIO 访问的函数均定义为驱动函数"的准则一致。

在识别出驱动函数列表后，模块进一步分析驱动函数之间的调用关系，构建函数调用图。调用图以驱动函数为节点，函数调用关系为边，记录调用方、被调用方、调用位置等信息。调用图记录了函数间的调用依赖关系，有助于理解函数的调用上下文和语义关联。

\begin{lstlisting}[style=CodeQL,language=CodeQL, caption={驱动函数调用关系查询},label=lst:driver-call-query]
from DriverToCall call
select call.toString2(),
       call.getTargetFunction().toString(),
       call.getEnclosingFunction().toString(),
       call.isDirectCall()
\end{lstlisting}

该查询从 \texttt{DriverToCall} 关系中提取驱动函数之间的调用关系，输出调用边的字符串表示、目标函数名、封装函数名以及调用是否为直接调用等信息，如代码~\ref{lst:driver-call-query}所示。

\subsection{数据模型实现}

为便于后续模块使用，静态分析结果被组织为统一的数据模型。数据模型采用 Python 数据类（dataclass）定义，支持序列化和反序列化，包括函数信息模型（\texttt{FunctionInfo}，包含函数名、签名、MMIO 访问列表等）、MMIO 访问点模型（\texttt{MmioAccessInfo}，包含结构体类型、字段名、访问类型等）和项目分析结果模型（包含调用图、驱动函数列表等），通过 JSON 序列化存储在缓存层供后续模块查询。

\subsubsection{查询工具函数实现}

各模块的数据通过一组工具函数对外暴露，供智能体在 LangGraph 工作流中调用。工具函数基于 LangChain 的 \texttt{@tool} 装饰器实现，将底层的数据查询封装为 LLM 可直接调用的工具接口。表~\ref{tab:tools}列出了系统中的主要工具函数及其功能。

\begin{table}[htbp]
\centering
\caption{主要工具函数}
\label{tab:tools}
\footnotesize
\setlength{\tabcolsep}{3pt}
\begin{tabular}{@{}lp{10.5cm}@{}}
\toprule
\textbf{工具函数} & \textbf{功能说明} \\
\midrule
GetFunctionInfo & 检索函数源码、签名、参数列表、返回类型、文件位置等信息 \\
GetMMIOFunctionInfo & 返回函数的所有 MMIO 访问点（结构体类型、字段名、访问类型、上下文代码） \\
GetFunctionCallStack & 提取函数的调用者列表和被调用者列表，返回调用关系双向视图 \\
GetStructOrEnumInfo & 查询数据结构或枚举类型的定义 \\
GetDriverInfo & 获取驱动模块的上下文信息 \\
\bottomrule
\end{tabular}
\end{table}

上述工具基于静态分析结果查询，内部实现了缓存机制避免重复访问。此外，系统还提供基于仿真运行结果的查询工具（如 MMIO 函数调用日志、函数调用栈信息），在仿真执行智能体中进行故障诊断时使用（见第~\ref{sec:hook-mechanism}~节）。各工具函数的返回结果均经过格式化处理，以结构化文本形式呈现给 LLM，避免原始 JSON 数据对模型推理的干扰。

\subsection{分析流程实现}

静态分析模块的分析流程包括四个步骤：（1）调用 \texttt{codeql database create} 构建包含 AST、控制流图和数据流图的 CodeQL 数据库；（2）依次执行前述 MMIO 表达式收集、常量基址数据流追踪、驱动函数判定和调用关系四组查询，将结果转换为 Python 对象；（3）将各查询结果进行关联聚合，构建完整的函数-MMIO 访问映射和函数调用图；（4）通过 JSON 序列化将聚合结果持久化到文件系统，支持后续模块快速查询和断点续跑。

\section{驱动函数分析、替换与编译模块实现}

\subsection{智能体工作流实现}

\subsubsection{驱动函数分析与替换代码生成模块实现}

分析智能体的设计见第~\ref{sec:analyzer-workflow}~节。在实现上，智能体基于 LangGraph 的 \texttt{StateGraph} 类构建，通过 \texttt{@tool} 装饰器将 GetFunctionInfo、GetFunctionCallStack、GetMMIOFunctionEmulateInfo 等工具函数注册到框架中，通过 \texttt{compile()} 方法编译为可执行的 \texttt{CompiledGraph} 对象，由执行器调度节点执行顺序并管理状态传递。

\subsubsection{替换代码修正与编译集成模块实现}
\label{sec:builder-fixer-impl}

编译修复智能体的设计见第~\ref{sec:builder-workflow}~节。

\textbf{代码集成机制}：系统采用"原位替换"策略将替换代码写入工程。对于每个需要替换的驱动函数，系统在目标源文件中保留原始函数定义的快照（用于回滚），然后将替换实现写入函数体。同时，系统在工程中添加弱符号桩函数文件，为仿真注入接口提供链接占位（详见第~\ref{sec:hook-mechanism}~节）。

\textbf{编译构建与错误修复流程}：编译修复智能体将替换代码应用到工程后执行编译构建，若编译失败则自动进入错误解析与修复循环。替换代码应用通过 ReplaceFuncs 工具实现，工程构建通过 BuildProject 工具执行 CMake 构建。内置的编译错误解析器支持 GCC、Clang、ARMCC 等多种编译器，能够将诊断信息转换为结构化错误对象。编译修复智能体按照设计采用"锁定首个错误"策略，优先修复第一个错误涉及的函数，修复策略由 LLM 根据错误信息和函数上下文智能选择。每次修复后重新触发编译，直至编译成功或达到最大修复轮次。

\subsubsection{替换代码验证机制实现}

第~\ref{chap:design}~章定义了替换代码的四层验证机制（语法正确性 $V_1$、签名匹配 $V_2$、返回值检查 $V_3$、替换规则检查 $V_4$），本节阐述各层验证的具体实现技术。

语法正确性检查通过 Python 的正则表达式和括号匹配算法实现，检查括号匹配、语句完整性、字符串闭合和函数体非空。签名匹配检查通过正则表达式从替换代码首行解析函数名、返回类型和参数列表，与 FunctionInfo 中的签名逐字段比对。返回值检查通过简化的控制流分析识别分支和循环出口，验证所有执行路径的返回值类型与声明兼容（如 \texttt{HAL\_StatusTypeDef} 的返回值须为该枚举的有效成员）。替换规则检查基于正则表达式模式匹配和关键字检索，执行策略特定的约束检查（如 Skip 策略不得包含 MMIO 寄存器访问，Redirect 策略须包含至少一个 \texttt{HAL\_BE\_*} 调用，Event 策略须保留回调函数调用，ReturnOK 策略的所有返回语句须返回成功状态值）。各层检查未通过时返回具体的违规描述，附加到分析智能体上下文中触发重新生成。

\subsection{替换策略实现}

各替换策略的设计意图与语义定义见第~\ref{sec:driver-replacement-strategy}~节，本节通过代码示例展示各策略的具体替换效果。代码示例均取自实际测试固件的替换结果，左侧为原始代码，右侧为替换后的代码。

\subsubsection{Skip 策略实现}

Skip 策略移除硬件交互操作但保留函数签名，适用于 INIT 和 LOOP 类函数。代码~\ref{lst:skip-strategy}分别展示了初始化类和轮询类函数的替换效果：初始化类函数移除寄存器配置后仅保留参数透传，轮询类函数移除基于硬件标志位的 while 循环后直接返回成功状态。

\begin{lstlisting}[style=CStyle,abovecaptionskip=0pt,belowcaptionskip=0pt,caption=Skip 策略替换示例（初始化类和轮询类）,label=lst:skip-strategy]
// 原始代码（初始化类）
void HAL_UART_Enable(UART_HandleTypeDef *huart) {
    huart->Instance->CR1 |= USART_CR1_UE;
}

// 替换代码（Skip 策略）
void HAL_UART_Enable(UART_HandleTypeDef *huart) {
    (void)huart;
}

// 原始代码（轮询类）
HAL_StatusTypeDef HAL_UART_WaitOnFlag(UART_HandleTypeDef *huart, uint32_t Flag, uint32_t Timeout) {
    while (!(__HAL_UART_GET_FLAG(huart, Flag))) {
        if (Timeout != HAL_MAX_DELAY) {
            if ((Timeout--) == 0) return HAL_TIMEOUT;
        }
    }
    return HAL_OK;
}

// 替换代码（Skip 策略）
HAL_StatusTypeDef HAL_UART_WaitOnFlag(UART_HandleTypeDef *huart, uint32_t Flag, uint32_t Timeout) {
    (void)huart;
    (void)Flag;
    (void)Timeout;
    return HAL_OK;
}
\end{lstlisting}

\subsubsection{Redirect 策略实现}

Redirect 策略将硬件寄存器轮询替换为仿真注入函数调用，同时保留非 I/O 业务逻辑。代码~\ref{lst:redirect-uart}展示了定长接收场景（UART），代码~\ref{lst:redirect-enet}展示了变长报文接收场景（以太网驱动）。

\begin{lstlisting}[style=CStyle,abovecaptionskip=0pt,belowcaptionskip=0pt,caption=Redirect 策略替换示例（UART 接收函数）,label=lst:redirect-uart]
// 原始代码
HAL_StatusTypeDef HAL_UART_Receive(UART_HandleTypeDef *huart, uint8_t *pData,
                                    uint16_t Size, uint32_t Timeout) {
    while (!(huart->Instance->SR & USART_SR_RXNE));
    *pData = huart->Instance->DR;
    return HAL_OK;
}

// 替换代码（Redirect 策略）
HAL_StatusTypeDef HAL_UART_Receive(UART_HandleTypeDef *huart, uint8_t *pData,
                                    uint16_t Size, uint32_t Timeout) {
    // 重定向到仿真输入函数
    HAL_BE_In(pData, Size);
    // 保留非 I/O 业务逻辑
    huart->RxXferCount = 0;
    huart->RxXferSize = Size;
    huart->RxState = HAL_UART_STATE_READY;
    return HAL_OK;
}
\end{lstlisting}

\begin{lstlisting}[style=CStyle,abovecaptionskip=0pt,belowcaptionskip=0pt,caption=Redirect 策略替换示例（变长报文接收）,label=lst:redirect-enet]
// 变长报文接收
curBuffDescrip = rxBdRing->rxBdBase + rxBdRing->rxGenIdx;
if (rxBuffer->buffer != NULL && rxBuffer->length > 0) {
    HAL_BE_ENET_ReadFrame(rxBuffer->buffer, rxBuffer->length);
}
rxBdRing->rxGenIdx = ENET_IncreaseIndex(rxBdRing->rxGenIdx, rxBdRing->rxRingLen);
\end{lstlisting}

\subsubsection{Event 策略实现}

Event 策略绕过硬件中断标志的判断条件，直接触发事件通知和回调函数。代码~\ref{lst:event-strategy}展示了 ETH 中断处理函数的替换效果。

\begin{lstlisting}[style=CStyle,abovecaptionskip=0pt,belowcaptionskip=0pt,caption=Event 策略替换示例,label=lst:event-strategy]
// 原始代码
void HAL_ETH_IRQHandler(ETH_HandleTypeDef *heth) {
    if (__HAL_ETH_DMA_GET_FLAG(heth, ETH_DMA_FLAG_R)) {
        HAL_ETH_RxCpltCallback(heth);
    }
    __HAL_ETH_DMA_CLEAR_FLAG(heth, ETH_DMA_FLAG_R);
}

// 替换代码（Event 策略）
void HAL_ETH_IRQHandler(ETH_HandleTypeDef *heth) {
    // 直接触发接收完成回调事件
    HAL_ETH_RxCpltCallback(heth);
    // 无需检查硬件标志或清除寄存器
}
\end{lstlisting}

\subsubsection{Steering 策略实现}

Steering 策略通过控制流改写将因硬件依赖而不可达的执行路径导向到仿真可达的分支，适用于 DMA+RingBuffer 等复杂状态机场景。代码~\ref{lst:steering-original}和代码~\ref{lst:steering-replaced}展示了 NXP 网卡驱动中 ENET\_GetRxFrame 函数的替换对比：原始代码通过轮询硬件 buffer descriptor 状态标志遍历接收数据（仿真中死循环），替换后跳过硬件标志扫描循环，在当前 buffer descriptor 处直接调用 HAL\_BE\_ENET\_ReadFrame 注入仿真数据，同时保留 Ring 索引更新和状态字段维护逻辑。

\begin{lstlisting}[style=CStyle,abovecaptionskip=0pt,belowcaptionskip=0pt,caption=Steering 策略原始代码示例（NXP ENET\_GetRxFrame）,label=lst:steering-original]
// 原始代码：依赖硬件标志扫描 Ring，仿真中将陷入死循环
index = rxBdRing->rxGenIdx;
while (0U == (curBuffDescrip->control & ENET_BUFFDESCRIPTOR_RX_EMPTY_MASK))
{
    if (0U != (curBuffDescrip->control & ENET_BUFFDESCRIPTOR_RX_LAST_MASK)) {
        result = ENET_GetRxFrameErr(curBuffDescrip, &rxFrame->rxFrameError);
        break;
    }
    index = ENET_IncreaseIndex(index, rxBdRing->rxRingLen);
    curBuffDescrip = rxBdRing->rxBdBase + index;
    if (index == rxBdRing->rxGenIdx) break;
}

// 后续：遍历 BD，拷贝数据到 rxFrame
do {
    newBuff = handle->rxBuffAlloc(base, handle->userData, ringId);
    if (newBuff != 0U) {
        rxBuffer = &rxFrame->rxBuffArray[index];
        if (0U != (curBuffDescrip->control & ENET_BUFFDESCRIPTOR_RX_LAST_MASK)) {
            isLastBuff = true;
            rxFrame->totLen = curBuffDescrip->length;
            rxBuffer->length = curBuffDescrip->length - buffLen;
        } else {
            rxBuffer->length = curBuffDescrip->length;
            buffLen += rxBuffer->length;
        }
        // ... buffer 管理、BD 状态更新 ...
    }
} while (!isLastBuff);
\end{lstlisting}

替换后的代码如代码~\ref{lst:steering-replaced}所示。

\begin{lstlisting}[style=CStyle,abovecaptionskip=0pt,belowcaptionskip=0pt,caption=Steering 策略替换代码示例（NXP ENET\_GetRxFrame）,label=lst:steering-replaced]
// 替换代码：跳过硬件标志扫描，直接注入数据
curBuffDescrip = rxBdRing->rxBdBase + rxBdRing->rxGenIdx;
index = 0;
do {
    newBuff = (uintptr_t)(uint8_t *)handle->rxBuffAlloc(base, handle->userData, ringId);
    if (newBuff != 0U) {
        rxBuffer = &rxFrame->rxBuffArray[index];
        rxBuffer->buffer = (void *)(uint8_t *)address;

        if (0U != (curBuffDescrip->control & ENET_BUFFDESCRIPTOR_RX_LAST_MASK)) {
            isLastBuff       = true;
            rxFrame->totLen  = curBuffDescrip->length;
            rxBuffer->length = curBuffDescrip->length - buffLen;
        } else {
            rxBuffer->length = curBuffDescrip->length;
            buffLen += rxBuffer->length;
        }
    }

    // Steering 核心：注入仿真数据
    if (rxBuffer->buffer != NULL && rxBuffer->length > 0)
        HAL_BE_ENET_ReadFrame(rxBuffer->buffer, rxBuffer->length);

    // 保留 Ring 更新，防止卡在同一 BD
    index++;
    rxBdRing->rxGenIdx = ENET_IncreaseIndex(rxBdRing->rxGenIdx, rxBdRing->rxRingLen);
    curBuffDescrip     = rxBdRing->rxBdBase + rxBdRing->rxGenIdx;
} while (!isLastBuff);
\end{lstlisting}

\subsubsection{ReturnOK 策略实现}

ReturnOK 策略适用于 PUREDRV 类中返回值为状态类型的函数。与 Skip 策略不同，ReturnOK 策略不生成替换代码，而是在仿真执行阶段通过 Unicorn 钩子直接拦截原始函数入口，根据 LLM 预先分析出的返回类型枚举值返回零值（如 HAL\_OK 的枚举值为 0）。这意味着被标记为 ReturnOK 的函数在 C 源码层面完全不修改，仅在仿真运行时被动态接管。代码~\ref{lst:returnok-strategy}以 GPIO 去初始化函数为例展示了 ReturnOK 策略的处理方式。

\begin{lstlisting}[style=CStyle,abovecaptionskip=0pt,belowcaptionskip=0pt,caption=ReturnOK 策略处理示例,label=lst:returnok-strategy]
// 原始函数（不生成替换代码，源码保持不变）
HAL_StatusTypeDef HAL_GPIO_DeInit(GPIO_TypeDef *GPIOx, uint32_t GPIO_Pin) {
    if (GPIOx == NULL) return HAL_ERROR;
    GPIOx->MODER &= ~(GPIO_MODER_MODE0 << (GPIO_Pin * 2));
    GPIOx->OTYPER &= ~(GPIO_OTYPER_OT0 << GPIO_Pin);
    GPIOx->OSPEEDR &= ~(GPIO_OSPEEDR_OSPEED0 << (GPIO_Pin * 2));
    GPIOx->PUPDR &= ~(GPIO_PUPDR_PUPD0 << (GPIO_Pin * 2));
    return HAL_OK;
}

// 仿真执行阶段：Unicorn hook 拦截该函数入口
// Python 处理器：读取返回类型 -> HAL_StatusTypeDef -> HAL_OK=0
// 直接向 r0 寄存器写入 0 并跳过函数体
\end{lstlisting}

\subsection{LLM 提示词工程实现}
\label{sec:prompt-engineering}

大语言模型的分类和生成质量直接影响替换代码的正确性，因此提示词的设计是系统的关键技术实现之一。本节阐述分析智能体和修复智能体的提示词模板结构设计。提示词工程的目标是：（1）为 LLM 提供函数上下文信息，使其能够理解函数语义；（2）通过结构化的分类规则指导 LLM 做出一致的分类决策；（3）通过生成约束限制 LLM 的输出格式，降低后处理复杂度。

\subsubsection{分类提示词模板设计}

分析智能体的提示词模板采用分阶段结构设计，包含角色定义、工具说明、分类策略和输出格式四个部分。其整体骨架如代码~\ref{lst:classifier-prompt-overview}所示。

\begin{lstlisting}[style=CStyle,abovecaptionskip=0pt,belowcaptionskip=0pt,caption={分类提示词模板骨架},label=lst:classifier-prompt-overview]
**Role**: You are an embedded software engineer tasked with
replacing driver library functions to eliminate dependencies
on peripheral hardware while preserving normal functionality
and MCU-related operations (including OS scheduling and
interrupt triggering).
**Overall Process**:
1. Information Gathering: Obtain the function implementation
   and context using available tools.
2. Analysis Phase: Analyze MMIO accesses, call relationships,
   data structures, and core logic.
3. Classification Phase: Categorize the function into one of
   the defined types.
4. Processing Phase:
   - RECV, IRQ, INIT, LOOP: Generate replacement code.
   - PUREDRV, NODRIVER: Classification only.
**Tools**: GetFunctionInfo, GetMMIOFunctionInfo,
  GetStructOrEnumInfo, GetFunctionCallStack, GetDriverInfo
**Categories**: RECV / IRQ / INIT / LOOP / PUREDRV /
  NODRIVER
**Rules**: Do not change return type; no new includes;
  preserve OS calls; only use HAL_BE_* helpers.
**Output Format** (JSON):
{
  "function_name": "<name>",
  "function_type": "<RECV|IRQ|INIT|LOOP|...>",
  "functionality": "<brief description>",
  "classification_reason": "<detailed explanation>",
  "has_replacement": <true|false>,
  "function_replacement": "<code or empty string>"
}
\end{lstlisting}

\textbf{角色定义与工具说明}。提示词以角色设定开场，将 LLM 定位为嵌入式软件工程师，明确任务是消除驱动函数的硬件依赖。随后声明五个可调用的工具函数及其调用顺序，如代码~\ref{lst:prompt-tools}所示。

\begin{lstlisting}[style=CStyle,abovecaptionskip=0pt,belowcaptionskip=0pt,caption={提示词中的工具声明与调用顺序},label=lst:prompt-tools]
1. GetFunctionInfo(func_name): Get function implementation.
   -> Always call this first.
2. GetMMIOFunctionInfo(func_name): Identify hardware accesses.
   -> Call immediately after GetFunctionInfo.
3. GetStructOrEnumInfo(struct_name): Query data structures.
4. GetFunctionCallStack(func_name): Get call relationships.
5. GetDriverInfo(driver_name): Get driver module context.
   Tool Sequence: GetFunctionInfo -> GetMMIOFunctionInfo
                 -> [3|4|5] (as needed)
\end{lstlisting}

提示词要求 LLM 先调用 GetFunctionInfo 获取函数源码，再调用 GetMMIOFunctionInfo 确认硬件访问方式，最后按需调用其余工具补充上下文。分类策略部分将七类函数的识别特征、替换策略和关键约束以结构化形式嵌入，按 RECV $>$ IRQ $>$ INIT $>$ LOOP 的优先级排列，减少高影响类别的误分类。代码生成约束限制 LLM 不得修改返回类型、不得引入新依赖、必须保留 OS 调用、仅使用 \texttt{HAL\_BE\_*} 注入函数，从源头减少幻觉问题。输出格式要求 JSON 结构化输出，便于后续自动解析。

编译修复智能体的提示词设计较为精炼，主要任务是接收编译器输出的错误信息和对应的函数名，要求 LLM 定位错误原因并生成修复后的替换代码，修复后重新触发编译验证。提示词约束修复时不得修改函数签名、不得引入新的头文件或外部声明，与分析智能体的通用替换规则一致。

\section{动态测试反馈闭环模块实现}

本节阐述动态测试反馈闭环模块的具体实现，其设计方案见第~\ref{chap:design}~章。

\subsection{仿真平台实现}

\subsubsection{Unicorn 引擎初始化与内存映射}

系统在初始化阶段通过 Unicorn 引擎的 C API 完成仿真环境的搭建。首先调用 \texttt{uc\_open(UC\_ARCH\_ARM, UC\_MODE\_THUMB, \&uc)} 创建 ARM Thumb 模式的模拟器实例，随后依次完成以下配置：（1）内存区域映射，根据目标芯片的链接脚本解析 Flash、RAM 和外设区域的起始地址与大小，调用 \texttt{uc\_mem\_map()} 为各区域分配独立的内存空间并配置读写权限，系统定义了 flash、ram、peripheral\_ram 等受保护的内存类型，防止不同区域被错误合并；（2）固件映像加载，支持 ELF 和 BIN 两种格式，通过自动识别代码段（而非数据段）作为 Flash 基址，调用 \texttt{uc\_mem\_write()} 将固件二进制写入虚拟内存；（3）寄存器初始化，通过 \texttt{uc\_reg\_write()} 设置程序计数器（PC）和栈指针（SP）等关键寄存器的初始值，确保固件从正确的入口点开始执行；（4）浮点单元支持，为兼容包含浮点运算的固件（如图像处理、网络协议栈），启用 VFP 浮点单元支持。

\subsubsection{钩子机制的实现}
\label{sec:hook-mechanism}

第~\ref{chap:design}~章阐述了本系统使用的四类钩子（代码钩子、基本块钩子、内存钩子、中断钩子）的设计原理，本节阐述其具体实现。

在钩子注册阶段，系统通过 \texttt{uc\_hook\_add()} 函数批量注册四类钩子，每类钩子绑定独立的 Python 回调函数。各回调函数接收 Unicorn 实例、触发地址和访问参数，将事件信息写入结构化的日志队列。具体实现中，代码钩子通过 \texttt{UC\_HOOK\_CODE} 类型注册，回调函数接收当前指令地址和指令大小，用于指令计数和执行轨迹记录；基本块钩子通过 \texttt{UC\_HOOK\_BLOCK} 类型注册，回调函数记录基本块入口地址，用于覆盖率统计；内存钩子通过 \texttt{UC\_HOOK\_MEM\_READ} 和 \texttt{UC\_HOOK\_MEM\_WRITE} 类型注册，回调函数接收访问地址、访问大小和数据内容，用于检测对未映射地址空间的非法访问；中断钩子通过 \texttt{UC\_HOOK\_INTR} 类型注册，回调函数接收中断号，用于处理 SVC 指令等软件中断。

在与驱动替换策略配合时，钩子机制承担了关键的拦截角色。系统在仿真初始化阶段通过 ELF 符号表解析所有 \texttt{HAL\_BE\_*} 注入函数的虚拟地址，然后在这些地址处注册代码钩子。当固件执行到注入函数入口时，钩子回调函数执行以下流程：（1）从 Unicorn 寄存器中读取函数参数（ARM 调用约定下参数通过 r0-r3 寄存器传递）；（2）调用宿主机侧的 Python handler 接口完成数据注入或输出操作；（3）将返回值写入 r0 寄存器；（4）从栈中读取返回地址并设置 PC 指针，跳过原函数体继续执行。通过这一机制，固件中调用 \texttt{HAL\_BE\_*} 系列接口时，执行流被拦截并转移到宿主机侧的处理例程。

\subsubsection{弱符号桩函数实现}

弱符号桩函数是连接固件代码和仿真环境的桥梁，其设计原理见第~\ref{chap:design}~章。代码~\ref{lst:weak-stubs}展示了桩函数的实现，所有函数均使用 \texttt{\_\_attribute\_\_((weak))} 修饰。在构建阶段作为链接占位符，在仿真运行阶段通过 Unicorn 钩子拦截并由宿主机侧 Python 处理器执行实际功能（详见第~\ref{sec:hook-mechanism}~节），实现同一份源码在硬件环境和仿真环境之间的零修改切换。

\begin{lstlisting}[style=CStyle,language=C,caption=弱符号桩函数实现示例,label=lst:weak-stubs]
__attribute__((weak)) void HAL_BE_In(uint8_t *buffer, uint16_t size) {}
__attribute__((weak)) void HAL_BE_Out(const uint8_t *data, uint16_t size) {}
__attribute__((weak)) void HAL_BE_ENET_ReadFrame(uint8_t *buffer,
                                                  uint16_t size) {}
__attribute__((weak)) uint32_t HAL_BE_return_0(void) { return 0; }
__attribute__((weak)) uint32_t HAL_BE_return_1(void) { return 1; }
__attribute__((weak)) void HAL_BE_Block_Read(uint8_t *buffer,
                                              uint32_t block_size,
                                              uint32_t block_count) {}
__attribute__((weak)) void HAL_BE_Block_Write(const uint8_t *buffer,
                                               uint32_t block_size,
                                               uint32_t block_count) {}
\end{lstlisting}

\subsubsection{运行日志记录模块实现}

运行日志记录模块通过前述的 Unicorn 钩子机制，将程序执行过程中的关键事件转化为结构化的日志记录。模块在每类钩子回调中将事件信息写入线程安全的日志队列，仿真结束后统一导出为 JSON 格式供故障诊断使用。日志记录包含以下四类信息：（1）指令级信息，利用代码钩子记录指令地址、指令大小和关键寄存器快照；（2）基本块级信息，利用基本块钩子记录每个基本块的入口地址和执行次数，据此构建代码覆盖率统计；（3）函数调用信息，在驱动函数入口和出口地址处设置代码钩子，记录调用地址、被调用函数名、参数值和返回值；（4）中断触发信息，利用中断钩子记录中断号、触发地址和处理器状态。这些日志信息在仿真运行过程中实时采集，为故障诊断和反馈闭环提供运行时数据基础。

\subsection{仿真执行智能体实现}

仿真执行智能体基于 LangGraph 框架实现，其设计见第~\ref{chap:design}~章。在实现层面，该智能体通过 \texttt{StateGraph} 类构建图式状态机，核心状态定义如下：

\begin{lstlisting}[style=Python,language=Python,caption=仿真执行智能体状态定义,label=lst:emulator-state]
@dataclass
class EmulatorState(TypedDict):
    messages: Annotated[list, add_messages]
    project_path: str
    build_config: dict
    replacement_functions: list
    iteration_count: int
    max_iterations: int
    processed_functions: list
    emulate_result: Optional[dict]
\end{lstlisting}

智能体的图结构包含以下节点：\texttt{emulate}（仿真运行节点，调用 EmulateProject 工具）、\texttt{diagnose}（故障诊断节点，调用 GetMMIOFunctionEmulateInfo 和 GetFunctionCallsEmulateInfo 工具）、\texttt{fix}（修复节点，调用 Fixer 子图）、\texttt{build}（构建节点，调用 Builder 子图）。节点间的条件转移逻辑根据仿真结果决定：仿真成功时直接终止；失败时进入诊断节点；诊断完成后进入修复节点；修复完成后进入构建节点；构建完成后返回仿真节点进行回归验证。每个节点的执行由 LLM 驱动决策，LLM 根据当前状态和工具返回的结果决定下一步操作。

\subsection{反馈闭环实现}

\subsubsection{多智能体协同的 LangGraph 实现}

第~\ref{chap:design}~章设计了由仿真执行智能体、修复智能体和构建智能体组成的多智能体架构，本节阐述其基于 LangGraph 的具体实现。

三个智能体通过 LangGraph 的子图（Subgraph）机制集成。仿真执行智能体作为主图，修复智能体和构建智能体分别编译为独立的 \texttt{CompiledGraph} 对象，通过工具调用接口被主图节点调用。这种设计使得各智能体拥有独立的状态空间和工具链，同时通过标准化的输入输出接口实现数据交换。图~\ref{fig:multi-agent-langgraph}展示了多智能体协同的 LangGraph 图结构。

\begin{figure}[htbp]
\centering
\begin{tikzpicture}[
    node distance=1.2cm and 1.8cm,
    mainbox/.style={draw, rounded corners=4pt, minimum width=2.8cm, minimum height=0.9cm, align=center, fill=blue!8, font=\small},
    subbox/.style={draw, rounded corners=4pt, minimum width=2.6cm, minimum height=0.9cm, align=center, fill=orange!10, font=\small},
    tool/.style={draw, rounded corners=2pt, minimum width=2.2cm, minimum height=0.6cm, align=center, fill=gray!12, font=\scriptsize},
    arr/.style={-{Stealth[length=5pt]}, thick},
    darr/.style={-{Stealth[length=5pt]}, thick, dashed},
]
\node[mainbox] (emulate) {Emulate\\仿真运行};
\node[mainbox, right=of emulate] (diagnose) {Diagnose\\故障诊断};
\node[mainbox, right=of diagnose] (fix) {Fix\\闭环修复};
\node[mainbox, right=of fix] (build) {Build\\闭环构建};

\node[tool, below=0.7cm of emulate] (t1) {EmulateProject};
\node[tool, below=0.7cm of diagnose] (t2) {\begin{tabular}{@{}c@{}}GetMMIOInfo\\GetCallStack\end{tabular}};
\node[tool, below=0.7cm of fix] (t3) {Fixer 子图};
\node[tool, below=0.7cm of build] (t4) {Builder 子图};

\draw[arr] (emulate) -- (diagnose);
\draw[arr] (diagnose) -- (fix);
\draw[arr] (fix) -- (build);
\draw[arr] (build.north) -- ++(0,0.6) -| node[above, pos=0.25, font=\scriptsize] {回归验证} (emulate.north);

\draw[darr] (emulate) -- (t1);
\draw[darr] (diagnose) -- (t2);
\draw[darr] (fix) -- (t3);
\draw[darr] (build) -- (t4);

\node[draw, dashed, rounded corners=6pt, fit=(emulate)(diagnose)(t1)(t2), inner sep=8pt, label={[font=\small]above:仿真执行智能体（主图）}] {};

\node[subbox, below=0.7cm of t3] (fixer) {闭环修复智能体};
\draw[darr] (t3) -- (fixer);

\node[subbox, below=0.7cm of t4] (builder) {闭环构建智能体};
\draw[darr] (t4) -- (builder);

\draw[arr, orange!70!black] (fixer) -- node[right, font=\scriptsize] {ReplacementUpdate} (builder);
\draw[arr, orange!70!black] (builder.west) -| node[below, pos=0.25, font=\scriptsize] {BuildOutput} (t4.south);

\node[left=0.3cm of emulate, font=\scriptsize, align=right] {success=false};
\node[above=1.2cm of fix, font=\small\bfseries] {};

\node[left=0.8cm of emulate, font=\small, align=center, text=green!50!black] {success=true\\$\to$ 终止};
\end{tikzpicture}
\caption{多智能体协同的 LangGraph 图结构}
\label{fig:multi-agent-langgraph}
\end{figure}

修复智能体接收仿真执行智能体传递的结构化输入，包括标准输出、标准错误、MMIO 函数调用日志和函数调用栈信息，输出为 \texttt{ReplacementUpdate} 对象（包含函数名、替换代码和修复原因）。构建智能体接收修复智能体的输出，执行代码应用和编译构建，输出为 \texttt{BuildOutput} 对象（包含退出码、编译输出和错误信息）。所有智能体间通信均通过预定义的数据结构进行，避免自由格式文本传递导致的理解偏差。

\subsubsection{闭环工作流的具体实现}

第~\ref{chap:design}~章给出了闭环协同工作流的核心逻辑，本节阐述其实现中的关键技术细节。图~\ref{fig:emulator-state-machine}展示了仿真执行智能体的状态机转移过程。

\begin{figure}[htbp]
\centering
\begin{tikzpicture}[
    node distance=1.4cm and 2.2cm,
    state/.style={draw, rounded corners=4pt, minimum width=2.4cm, minimum height=0.8cm, align=center, fill=blue!8, font=\small},
    decision/.style={draw, diamond, aspect=2.2, inner sep=1pt, align=center, fill=yellow!12, font=\small},
    terminal/.style={draw, rounded corners=8pt, minimum width=2cm, minimum height=0.7cm, align=center, font=\small},
    arr/.style={-{Stealth[length=5pt]}, thick},
]
\node[state] (init) {环境准备\\加载固件};
\node[state, right=of init] (run) {仿真运行};
\node[decision, right=of run] (check) {成功?};
\node[terminal, right=1.5cm of check, fill=green!12] (done) {终止\\输出结果};
\node[state, below=of check] (diag) {故障诊断\\PC分析+日志};
\node[state, left=of diag] (repair) {修复执行\\锁定单函数};
\node[state, left=of repair] (rebuild) {重新构建\\交叉编译};
\node[terminal, below=0.8cm of diag, fill=red!10] (fail) {终止\\失败报告};

\draw[arr] (init) -- (run);
\draw[arr] (run) -- (check);
\draw[arr] (check) -- node[above, font=\scriptsize] {是} (done);
\draw[arr] (check) -- node[right, font=\scriptsize] {否} (diag);
\draw[arr] (diag) -- (repair);
\draw[arr] (repair) -- (rebuild);
\draw[arr] (rebuild.north) -- ++(0,0.5) -| node[above, pos=0.25, font=\scriptsize] {回归验证} (run.south);

\draw[arr] (diag) -- node[right, font=\scriptsize] {达到上限/停滞} (fail);
\end{tikzpicture}
\caption{仿真执行智能体状态机转移图}
\label{fig:emulator-state-machine}
\end{figure}

在故障诊断环节，系统采用第~\ref{chap:design}~章设计的 PC 分析方法，将崩溃地址与静态分析输出的函数地址范围匹配以定位问题函数。对于超时故障，系统分析 PC 分布直方图识别循环结构。

在修复执行环节，修复智能体按照设计采用"锁定单函数"策略，每轮仅修改一个函数。修复策略包括：为 LOOP 类函数添加超时退出条件、补充遗漏的状态字段更新、将不适配的 Skip 策略调整为 Redirect 策略等。

在回归验证环节，系统在每次构建成功后自动重新生成仿真配置并执行完整的仿真运行。最大迭代次数 $K$ 见第~\ref{chap:design}~章设计，达到限制后终止闭环并生成详细故障报告。

\subsubsection{终止条件与错误处理的实现}

闭环的终止判断通过 LangGraph 的条件边（conditional edge）实现。在仿真节点之后，条件函数检查仿真结果：若 \texttt{emulate\_result.success == true}，转移到终止节点；否则检查迭代计数是否达到上限、诊断是否能定位问题函数，决定继续迭代或终止。

系统实现了三层错误处理机制：（1）编译错误处理，由构建智能体在编译失败时触发，错误信息通过正则表达式解析后传递给修复智能体；（2）运行时错误处理，包括内存访问错误（UC\_ERR\_FETCH\_UNMAPPED、UC\_ERR\_WRITE\_UNMAPPED）、非法指令执行（UC\_ERR\_INSTRUCTION\_INVALID）和超时，错误信息从 Unicorn 异常中提取后传递给故障诊断模块；（3）框架级错误处理，当 LangGraph 达到递归深度限制（GraphRecursionError）时，系统捕获异常并分析对话历史，生成故障报告后退出。系统记录每次迭代的完整信息（失败现象、诊断结果、修复方案、验证结果），便于问题追溯和系统优化。

\subsubsection{修复智能体提示词模板设计}

闭环修复智能体的提示词与分析智能体的分类提示词不同，将错误诊断和修复方案生成整合在一个推理流程中，其骨架如代码~\ref{lst:fixer-prompt-overview}所示。

\begin{lstlisting}[style=CStyle,abovecaptionskip=0pt,belowcaptionskip=0pt,caption={修复提示词模板骨架},label=lst:fixer-prompt-overview]
**Role**: Embedded software engineer specializing in HAL
development and debugging. Your task is to actively and
independently analyze emulator error logs to identify
problematic functions and fix them.
**CRITICAL**: You must NOT ask for parameters. You MUST
first actively retrieve error logs using available tools.
**Workflow**:
Step 0: Retrieve emulator logs (function_calls_emulate_info,
        mmio_function_emulate_info)
Step 1: Analyze logs to identify problematic functions
Step 2: Get function context (GetReplaceFuncDetailsByFile)
Step 3: Generate fix and apply (UpdateFunctionReplacement)
**Constraints**: Preserve function signature; no new types
  or headers; preserve OS core functions (scheduler, context
  switch); do not stub VTOR/NVIC/SysTick writes.
\end{lstlisting}

修复提示词的设计要点包括：（1）主动日志获取——要求 LLM 首先调用 \texttt{function\_calls\_emulate\_info()} 和 \texttt{mmio\_function\_emulate\_info()} 获取完整的执行流和错误信息，从日志中自行定位问题函数，而非等待外部指定；（2）上下文完整获取——要求 LLM 在生成修复方案前必须调用 \texttt{GetReplaceFuncDetailsByFile} 获取该函数的历次替换历史，避免重复已失败的修复方案；（3）OS 核心函数保护——禁止 LLM 修改操作系统核心函数（如 \texttt{vTaskStartScheduler}、\texttt{rt\_schedule} 等），若崩溃发生在 OS 核心函数中则应查找附近的驱动初始化代码问题；（4）步数预算控制——限制 LLM 对最多 3--5 个可疑函数进行分析后立即生成修复方案，避免无限调用分析工具耗尽步数。

\subsubsection{仿真执行智能体与闭环构建智能体提示词设计}

仿真执行智能体作为动态测试反馈闭环的主控调度单元，其提示词定义了闭环的验证标准和状态驱动的调试工作流，如代码~\ref{lst:emulator-prompt}所示。提示词的设计特点是：以"无死循环、无异常、到达 main 函数"三项条件作为成功判定依据；严格规定 Emulate $\to$ Diagnose $\to$ Build 的三步循环且不允许跳步；明确告知 LLM 修复是否成功只能通过仿真运行结果判断，而非通过静态分析推断。闭环构建智能体的提示词与分析智能体中的编译修复智能体类似，额外要求 LLM 在遇到大量编译错误时逐函数修复，避免上下文溢出。

\begin{lstlisting}[style=CStyle,abovecaptionskip=0pt,belowcaptionskip=0pt,caption={仿真执行智能体提示词骨架},label=lst:emulator-prompt]
**Task**: Iteratively debug and validate the firmware after
replacing hardware-coupled functions with mock stubs.
**Validation Criteria** (all must be true):
1. Program does NOT enter an infinite loop.
2. No exceptions (illegal memory access, bus faults).
3. Execution successfully reaches the main() function.
**Debugging Workflow** (strict state-driven cycle):
1. Emulate -> Check success field.
   If success: true -> STOP (task complete).
   If success: false -> Proceed to step 2.
2. Diagnose (exactly once per failure cycle)
   -> Call Fixer to analyze logs and produce code fixes.
   -> After receiving Fixer response -> Proceed to step 3.
3. Build -> Recompile with proposed changes.
   If build fails -> Go back to step 2.
   If build succeeds -> Go to step 1 (re-emulate).
**Critical Rules**:
- One Fixer call per failure iteration.
- Always call Builder after Fixer.
- Validation is empirical: only Emulator's success field
  confirms a fix, not further analysis.
\end{lstlisting}

闭环构建智能体的提示词与分析智能体中的编译修复智能体类似，主要负责应用修复代码并执行编译构建，额外要求 LLM 在遇到大量编译错误时逐函数修复。

\section{本章小结}

本章围绕第三章的设计方案，采用 CodeQL、LangGraph、Unicorn 等技术栈完成了系统的具体实现：基于 CodeQL 实现了三种 MMIO 定位模式的查询规则和数据流追踪；基于 LangGraph 实现了分析、编译修复、仿真执行和闭环修复四个智能体的工作流；基于 Unicorn 构建了仿真平台，通过弱符号桩函数和 Python 处理器实现固件与仿真环境的数据交互。实现结果表明，系统能够在无人干预的情况下完成从静态分析到动态反馈的完整流程。
